% source: RKF/theorum/29_morphic_smriti_stage15_plan.md
\hypertarget{stage-1.5-morphic-identity-recognition-completion-and-smriti-tails}{%
\chapter{Stage 1.5: Morphic Identity, Recognition Completion, and Smriti Tails}\label{stage-1.5-morphic-identity-recognition-completion-and-smriti-tails}}

\hypertarget{stage-rule}{%
\section{1. Stage rule}\label{stage-rule}}

Stage 1.5 is a writing and theorem-design stage. It does not claim a test pass.
The generalized calculus is written first; exact tests are added only after the
current Stage-1 packet has been run and reviewed.

This preserves the project order:

\begin{verbatim}
native definition
-> theorem
-> exact generalized example
-> fail-closed tests
-> domain specialization
-> numerical certification.
\end{verbatim}

\begin{center}\rule{0.5\linewidth}{0.5pt}\end{center}

\hypertarget{buried-material-retained}{%
\section{2. Buried material retained}\label{buried-material-retained}}

The following ideas are retained from the buried calculus chapters:

\begin{verbatim}
objects as stabilized morphic lineages;
trans-cut morphisms;
Recognition-Seam product and chain residues;
recovered middle identity;
Hilbert shadow versus recognition completion;
Madhava finite recursion plus Smriti tail;
projection leakage and recognition-null quotient;
circular/hyperbolic charts as optional representations.
\end{verbatim}

The following remain illustrative until represented rigorously:

\begin{verbatim}
specific Tamil morphic commutator constants;
physical interpretations of Crown-sector terms;
formal recognition-index direct sums without target groups;
curvature flux formulas without declared connections and domains.
\end{verbatim}

\begin{center}\rule{0.5\linewidth}{0.5pt}\end{center}

\hypertarget{generalized-example-e-recovered-middle-identity}{%
\section{3. Generalized example E: recovered middle identity}\label{generalized-example-e-recovered-middle-identity}}

Construct two arrows

\[
A\xrightarrow{\gamma}B_1,
\qquad
B_2\xrightarrow{\delta}C
\]

with identical shadow middle coordinates but different cut memories.

The example must verify:

\begin{verbatim}
shadow middle equality                         true
recovered middle identity                      false
naive matrix composition                       defined
RSC composition                                rejected
unledgered composition residue                 nonzero
\end{verbatim}

Then add an exact memory alignment arrow and verify that recovered identity and
composition are restored.

Purpose: prevent finite packets with the same visible coordinates but different
source/cut lineage from being spliced together.

\begin{center}\rule{0.5\linewidth}{0.5pt}\end{center}

\hypertarget{generalized-example-f-shadow-cauchy-but-not-recognition-cauchy}{%
\section{4. Generalized example F: shadow-Cauchy but not recognition-Cauchy}\label{generalized-example-f-shadow-cauchy-but-not-recognition-cauchy}}

Use a constant shadow vector

\[
x_n=x
\]

with a drifting integer memory

\[
k_n=n.
\]

Verify:

\begin{verbatim}
shadow norm distance                           0
shadow sequence Cauchy                         true
cut-memory sequence Cauchy                     false
recognition sequence Cauchy                    false
\end{verbatim}

Then use a lawfully ledgered memory quotient and verify recognition-Cauchy
closure in the quotient.

Purpose: demonstrate why norm convergence alone cannot close an
infinite-dimensional cut theorem.

\begin{center}\rule{0.5\linewidth}{0.5pt}\end{center}

\hypertarget{generalized-example-g-madhavasmriti-completion}{%
\section{5. Generalized example G: Madhava--Smriti completion}\label{generalized-example-g-madhavasmriti-completion}}

Choose an exact convergent operator series with finite partial recursion \(S_n\),
a predeclared correction \(C_n\), recognized limit \(B_\infty\), and exact tail
\(\mathsf S_n\):

\[
B_\infty=S_n+C_n+\mathsf S_n.
\]

The example must verify:

\begin{verbatim}
finite recursion generated independently       true
correction declared before evaluation           true
tail identity exact                             true
tail norm decreases                             true
retrospective fitted correction rejected        true
\end{verbatim}

Purpose: make the finite-to-infinite theorem use lawful tail memory rather than
visual convergence.

\begin{center}\rule{0.5\linewidth}{0.5pt}\end{center}

\hypertarget{generalized-example-h-morphic-object-stabilization}{%
\section{6. Generalized example H: morphic object stabilization}\label{generalized-example-h-morphic-object-stabilization}}

Construct a finite arrow family whose different admissible paths produce one
recovered object after ledger closure.

Verify:

\begin{verbatim}
local arrow residues                            zero
path shadows agree                              true
path cut memories agree or are ledger-equivalent true
stabilized recognition class                    unique
\end{verbatim}

Add a negative control in which endpoint shadows agree but path memory differs,
so the alleged object is not stabilized.

Purpose: implement the doctrine that objects are stabilized morphic lineages,
not primitive labels.

\begin{center}\rule{0.5\linewidth}{0.5pt}\end{center}

\hypertarget{generalized-example-i-recognition-complete-five-memory-limit}{%
\section{7. Generalized example I: recognition-complete five-memory limit}\label{generalized-example-i-recognition-complete-five-memory-limit}}

Build a directed exact-rational sequence

\[
Z_n=Y_n+M_n
\]

with:

\begin{verbatim}
recognized channel Cauchy                       true
memory channel Cauchy                           true
faithfulness residual                           -> 0
target dimension                                5
uniform reference floor                         positive
finite seam matrices                            converge
\end{verbatim}

Verify the quantitative bounds in the Recognition-Complete Finite-to-Infinite
Cut Theorem and show that

\[
\lambda_{\max}(B_n)+e_n<1
\]

certifies the limit form.

Negative controls:

\begin{verbatim}
shadow convergence with drifting memory;
rank-five numerical shadow without faithfulness;
positive determinant with two super-threshold modes;
changing the five-label orientation between refinements;
post-hoc correction fitted to the terminal residue.
\end{verbatim}

\begin{center}\rule{0.5\linewidth}{0.5pt}\end{center}

\hypertarget{test-order-after-writing}{%
\section{8. Test order after writing}\label{test-order-after-writing}}

After the user reports the Stage-1 output, implement tests in this order:

\begin{verbatim}
T1 transition differential and composition residues
T2 recovered middle identity
T3 recognition-Cauchy versus shadow-Cauchy
T4 Madhava--Smriti exact tail
T5 morphic stabilization
T6 target-faithful limit
T7 finite seam-matrix convergence
T8 adversarial negative controls
\end{verbatim}

No RH data enters these tests.

\begin{center}\rule{0.5\linewidth}{0.5pt}\end{center}

\hypertarget{rh-stage-after-generalized-validation}{%
\section{9. RH stage after generalized validation}\label{rh-stage-after-generalized-validation}}

Only after Stage 1.5 passes do we construct:

\[
Z_n^{\rm Weil}
=Z_n^{\rm prime}
\oplus Z_n^{\Gamma}
\oplus Z_n^{\partial}
\oplus Z_n^{\rm comp}.
\]

The first RH task is not an eigenvalue calculation. It is to prove recovered
identity and Smriti-complete convergence of these four event families on one
common cut carrier.

\begin{center}\rule{0.5\linewidth}{0.5pt}\end{center}

\hypertarget{claim-boundary}{%
\section{10. Claim boundary}\label{claim-boundary}}

\begin{verbatim}
Stage-1.5 theorem architecture                     WRITTEN
exact examples E--I                                PLANNED
fail-closed tests                                  DEFERRED UNTIL STAGE-1 REVIEW
RH specialization                                  DEFERRED
RH numerical certificate                           DEFERRED
\end{verbatim}

\begin{flushright}\small\texttt{RKF/theorum/29\_morphic\_smriti\_stage15\_plan.md}\end{flushright}
